\subsection{Point of Intersection of Lines} 
If we have two different lines that are not parallel, they \makeblank[1in] at exactly one point.
\begin{example}
Graph the lines described below, and find their point of intersection.
\[
\begin{system}
y=-\frac{2}{3}x+3\\
y=x-2
\end{system}
\]
\begin{icpic}
\setgraphsquare{5}
\GraphLarge
\end{icpic}
The lines intersect at \blankpair
\end{example}
Notice:
\begin{itemize}
\item $\makeblanksmall=-\frac{2}{3}(\makeblanksmall)+3$, so \blankpair is a solution to the first equation.
\item $\makeblanksmall=\makeblanksmall-2$, so \blankpair is a solution to the second equation.
\end{itemize}
This is called a \emph{system of equations}.
\begin{definition}[System of Equations]
A \term{system of equations} is a set of two or more equations that are to be solved simultaneously, that is, equations that should all be true at the same time.
\end{definition}
If we \makeblank[1in] a system of equations, the solution to the system is the point where the lines \makeblank[1in].